% Generated by eval/scripts/write_collinearity_explanation.py.
The practical question is whether the regression can separate the distance association from the specialization association. It includes distance, SPEC, model identity, and attack language. Some of these variables may contain overlapping information. For example, if SPEC is consistently higher for certain models, differences associated with SPEC also coincide with differences between those models. The model then has less information with which to distinguish the SPEC association from their baseline differences. This overlap is called \emph{collinearity}.

We check that overlap using the \emph{variance inflation factor} (VIF). A value of one indicates no linear overlap between a predictor and the other predictors; larger values indicate greater overlap.\footnote{VIF is calculated as $1/(1-R^2)$ by regressing the predictor on the other fitted design columns. This auxiliary $R^2$ concerns the predictor, not StrongREJECT scores. VIF is a design diagnostic; it does not directly multiply the robust GEE standard errors.} Here, ``other predictors'' includes the model and language indicators, as well as the other substantive predictor.

Table~\ref{tab:graded-diagnostics} reports distance VIF 2.46 and SPEC VIF 13.51 in the pooled fit. Distance has comparatively less overlap; SPEC has substantial overlap, making its separate adjusted contribution harder to isolate. This does not imply that SPEC is unimportant. Because its VIF exceeds the chosen threshold of five, we check simpler fits that include only one substantive predictor at a time.

The condition number checks the same general problem across the whole design, rather than for one predictor. Larger values indicate combinations of columns that are harder to distinguish numerically. It is 16.7 in the pooled fit and 27.2 when IF and CONS replace SPEC. These values depend on scaling and coding; our design includes cohort-standardized continuous predictors, an intercept, and treatment-coded model/language indicators. We report them alongside VIF without imposing a separate condition-number cutoff.
